\scanpage{51}
% @ line 1 / text
Considérons l'hom
% @ lines 2–3 / display
\[
 Q\longrightarrow\check{\mathbb P}(\unsure{\mathfrak g|g})
 =\underset{\substack{\shortparallel\\\Sigma}}{\mathbb P(\widetilde{\mathfrak g})}
 \quad(=\ \text{Schéma des ss-Algèbres spéciales de }M
 \tag*{(19)}
\]
% @ lines 4–6 / text
C'est évidemment un \uline{épi}, et si $A_\Sigma$ est ``la sous-algèbre spéc. universelle'' de $M$ -- ou \unsure{plutôt}, de $M_\Sigma$ -- on trouve donc un iso. can.
% @ line 7 / display
\[
 Q\simeq\operatorname{Spec}(A_\Sigma)\tag*{(20)}
\]
% @ left below (20) / text
\noindent\shortstack{\rule{1.7cm}{0.4pt}\\[-0.55em]\rule{1.7cm}{0.4pt}}

% @ lines 8–15 / text
Ceci étant, $A$ étant fixée, et $\underline L,\underline M$ étant des sous-fibrés inv. de $\widetilde{\mathfrak g}$ correspondants à des sections de $X$ (i.e. sur lesquels $q$ s'annule) à quelle condition le $\underline J$ \deleted{\unsure{que}}\inserted{$=L^\perp\otimes M$} $\subset M$ section de $Q$ qui leur correspond est-elle $\subset A$ ? (Les couples $\underline L,\underline M$ en question correspondent donc exactement aux sections de $Y=\operatorname{Spec} A$ sur $S$.)\unsure{.}

% @ line 16, first cancelled line / deletion
\deleted{Je dis qu'il \unsure{faut s.} que $L^{\unsure{\epsilon|{}}}\subset A^\perp$}
% @ lines 17–19, crossed proposition / text
\deleted{\uline{Prop} Conditions équivalentes, pour $A,L,M$ donnés} \deleted{\unsure{(i.e.}} \deleted{\unsure{$\Leftrightarrow L\subset$}}
% @ lines 18–19, bracketed hypotheses / display
\[
 \deleted{\unsure{\Leftrightarrow}\left\{\begin{aligned}
 &A\ \text{ss-Alg. spéciale de }M\\
 &L,M\ \text{ss-fibrés }\inserted{\text{inv de rang }1\text{ de }V}\ 
 \deleted{\text{satisfaisant }\operatorname{Tr},\det|L,M=0}
 \end{aligned}\right.}
\]
% @ lines 20–21, paired conditions / display
\[
 \deleted{\left.\begin{aligned}
 a)\quad &\underline L\subset A^\perp\quad(\text{i.e }\operatorname{Tr}\underline LJ=0\quad\text{si } A=\mathcal O_S\oplus J)\\
 b)\quad &\underline M\subset A^\perp\quad(\text{i.e. }\operatorname{Tr}\underline MJ=0\quad\text{————})
 \end{aligned}\right\}}
\]
% @ line 22 / display
\[
 \deleted{c)\quad L^\perp\otimes M\subset A}
\]
% @ lines 23–25, end at lower edge / text
\deleted{De plus, pour tt $L$\inserted{$\in X(S)$} \deleted{satisfaisant \unsure{a)}}, \inserted{$L\subset A^\perp$,} il existe un et un seul $M\in X(S)$ satisfaisant c) [ donc}
% archive sheet number: 50
